\section{Annexes}

\subsection{Hypothèses et dépendances}
\label{sec:hypotheses}

\begin{description}
  \item[Tado reste l'unique source.] L'application ne mesure ni ne commande rien par
    elle-même~: elle montre ce que Tado lui dit et lui transmet les ordres de l'utilisateur.
    Ce que Tado ne met pas à la disposition des applications tierces est abandonné.
  \item[Eau chaude et climatisation exclues.] La maison n'en a pas~; l'application ne les
    montre pas.
  \item[Présence par géolocalisation.] Le mode \emph{Auto} de Tado se fonde sur la position
    des appareils mobiles suivis, pas sur leur connexion au wifi~; ce choix est retenu. L'utilisateur n'a pas l'abonnement
    Auto-Assist, mais la géolocalisation fonctionne dans son application Android~; l'API
    accepte de rendre la présence à ce mode sans l'abonnement (vérifié le 28 octobre 2026).
  \item[Fonctions conditionnées à l'API Tado.] Le bien-être ambiant et les alertes
    reprennent ce que montre l'application Android Tado. Toute information que l'API de
    Tado ne fournit pas est abandonnée, avec la fonction qui en dépend.
  \item[Montre associée à un téléphone.] Une montre Pebble n'atteint le réseau que par le
    téléphone auquel elle est associée~; lorsqu'elle s'en éloigne, elle ne reçoit plus rien
    de Tado (\refexig{TPB-ACT-006}).
  \item[Programmation faite ailleurs.] Les plages horaires se règlent dans l'application
    mobile Tado~; l'application les lit, ne les écrit pas.
\end{description}

\subsection{Exemple suivi~: la maison d'essai}
\label{sec:maison-essai}

La maison ci-dessous illustre les exigences et sert de jeu d'essai
(§\ref{sec:verification}). Elle compte six pièces, un soir d'hiver, la maison étant
\emph{Présent} en mode \emph{Auto} (deux appareils à la maison).

\begin{table}[H]
  \centering\small
  \begin{tabularx}{\linewidth}{@{}lrrrlX@{}}
    \toprule
    \textbf{Pièce} & \textbf{Mesure} & \textbf{Demandée} & \textbf{Puissance} & \textbf{Mode} & \textbf{Couleur (§\ref{sec:affichage})} \\
    \midrule
    Salon         & 19,4~°C & 21~°C & 60~\% & Programme       & cyan \\
    Cuisine       & 20,1~°C & 19~°C & 0~\%  & Programme       & ambre \\
    Ch. parents   & 17,8~°C & 17~°C & 0~\%  & Programme       & bleu \\
    Ch. enfant    & 18,2~°C & 19~°C & 35~\% & Forcé, encore 1~h & cyan \\
    Bureau        & 19,9~°C & 20~°C & 15~\% & Programme       & ambre \\
    Salle de bain & 21,3~°C & 22~°C & 0~\%  & Programme       & orange \\
    \bottomrule
  \end{tabularx}
  \caption{Les six pièces de la maison d'essai}
  \label{tab:maison-essai}
\end{table}

Alertes en cours~: fenêtre ouverte au salon (chauffage coupé depuis 4~min)~; pile faible dans
la chambre de l'enfant (tête thermostatique).

Ce que l'exemple met à l'épreuve~:
\begin{itemize}
  \item le salon chauffe à 60~\% et une chambre à 35~\% (\refexig{TPB-PIE-002},
        \refexig{TPB-PIE-003})~;
  \item la cuisine est au-dessus de sa température demandée~: elle ne chauffe pas
        (\refexig{TPB-PIE-004})~;
  \item la chambre de l'enfant est forcée (\refexig{TPB-PIE-005}, \refexig{TPB-DET-002})~;
  \item les quatre tranches de couleur apparaissent (\refexig{TPB-AFF-001})~;
  \item la salle de bain, au-dessus de 21~°C, est la seule pièce orange~;
  \item deux alertes de natures différentes (\refexig{TPB-ALE-002}, \refexig{TPB-ALE-003})~;
  \item la maison est \emph{Présent} en mode \emph{Auto}, avec deux appareils
        (\refexig{TPB-PRE-001}, \refexig{TPB-PRE-003}).
\end{itemize}

\subsection{Écarts avec l'existant}
\label{sec:ecarts}

Le tableau ci-dessous, généré à partir des exigences, récapitule celles que l'application
ne satisfait pas aujourd'hui. L'application n'étant pas encore écrite, il reste vide.

\listeecarts

\subsection{Points ouverts}
\label{sec:points-ouverts}

Les douze points ouverts de la première rédaction ont été tranchés le 28 octobre 2026~;
leurs décisions figurent dans les exigences concernées. Aucun point ne reste ouvert.
% Réactiver \listepointsouverts dès qu'un nouveau \pointouvert apparaît.
